\subsection{پیشینه و مرور ادبیات}

پیشینه‌ی این پژوهش سه انتخاب روش‌شناختی را پوشش می‌دهد؛ استخراج کمیت‌های هندسی از عروق، استفاده از شبکه‌ی تصویری به‌عنوان خط پایه، و جداسازی بیماران یا معاینات هنگام تقسیم داده. تفاوت جامعه، دستگاه تصویربرداری و تعریف خروجی اجازه نمی‌دهد عددهای مطالعات مختلف بدون بررسی این جزئیات با نتایج حاضر مقایسه شوند.

\subsubsection{رتینوپاتی نارس و غربالگری \lr{Plus}}

رتینوپاتی نارس یا \lr{ROP}\enfoot{بیماری رشد غیرطبیعی عروق شبکیه در نوزاد نارس؛ معادل انگلیسی: \lr{Retinopathy of Prematurity}} در نوزادان با وزن تولد پایین شایع است و یکی از علل مهم نابینایی قابل پیشگیری در کودکان به‌شمار می‌رود. طبقه‌بندی بین‌المللی ویرایش سوم \lr{Plus}\enfoot{گشادشدگی و پیچ‌خوردگی پیش‌رونده‌ی عروق قطب خلفی؛ معادل انگلیسی: \lr{Plus}} و \lr{Pre-Plus}\enfoot{وضعیت میانی میان ظاهر طبیعی و پلاس؛ معادل انگلیسی: \lr{Pre-Plus}} را بخشی از یک طیف تغییرات عروقی می‌داند و بر ارزیابی عروق قطب خلفی تأکید می‌کند \cite{chiang2021}. تعریف \lr{Plus} بر گشادشدگی و پیچ‌خوردگی همین عروق استوار است و مرزهای آن با قضاوت متخصص تعیین می‌شود؛ توافق متخصصان بر این تشخیص کامل نیست. غربالگری به‌موقع این نوزادان بخشی از مراقبت استاندارد است \cite{fierson2018} و خودکارسازی از روی فوندوس می‌تواند دسترسی را بهبود دهد، اما ارزیابی منصفانه بدون جداسازی گروهی داده شرط پذیرش علمی است و عدد بالا روی تقسیم نادرست ارزش بالینی ندارد.

\subsubsection{سه خانواده روش}

شبکه‌ی کانولوشنی سرتاسری\enfoot{شبکه‌ای که از پیکسل خام تا برچسب نهایی یکپارچه آموزش می‌بیند؛ معادل انگلیسی: \lr{end-to-end CNN}} الگو را مستقیماً از رنگ و بافت تصویر می‌آموزد. مدل جدولی بر دانش قبلی درباره‌ی شکل رگ تکیه دارد و ویژگی‌های قابل‌نام‌گذاری تولید می‌کند، اما خطای ماسک عروق را نیز به طبقه‌بند منتقل می‌کند. ترکیب این دو مسیر فقط زمانی سودمند است که اطلاعات افزوده مکمل بازنمایی تصویر باشد؛ به همین دلیل نتیجه‌ی شاخه‌ی ترکیبی بدون مقایسه با بازنمایی تصویری تنها قابل تفسیر نیست.

\subsubsection{شبکه \lr{EfficientNet} و کتابخانه \lr{timm}}

شبکه \lr{EfficientNet-B5}\enfoot{معماری تصویری با مقیاس‌گذاری هم‌زمان عمق، عرض و وضوح؛ معادل انگلیسی: \lr{EfficientNet-B5}} با مقیاس‌گذاری ترکیبی عمق، عرض و وضوح، خط پایه‌ی تصویری این پروژه است \cite{tan2019} و پیاده‌سازی آن از کتابخانه‌ی \lr{timm}\enfoot{کتابخانه‌ی مدل‌های تصویری ازپیش‌آموزش‌دیده برای \lr{PyTorch}؛ معادل انگلیسی: \lr{timm}} می‌آید. آموزش با آنتروپی متقاطع وزن‌دار کلاس انجام شد و انتخاب دور آموزش روی \lr{AUC} اعتبارسنجی صورت گرفت تا آزمون در تصمیم‌ها وارد نشود.

\subsubsection{کتابخانه اندازه‌گیری عروقی}

کتابخانه‌ی \lr{PVBM}\enfoot{جعبه ابزار پایتون برای استخراج ویژگی‌های هندسی عروق از ماسک؛ معادل انگلیسی: \lr{PVBM}} ستون‌های ویژگی هندسی و چگالی را از ماسک دودویی عروق استخراج می‌کند \cite{fhima2022}. این پژوهش از همان خانواده‌ی اندازه‌گیری استفاده می‌کند، اما مجموعه‌ی نهایی آن به پنج نشانگر محدود شد. جداساز عروقی این پروژه \lr{MAnet}\enfoot{شبکه‌ی جداسازی معنایی با توجه چندمقیاسه؛ معادل انگلیسی: \lr{MAnet}} با رمزگذار \lr{ResNet34} است. ماسک مرجع متخصص برای داده‌ی این پروژه در دسترس نبود، پس صحت جداسازی در برابر مرز انسانی سنجیده نشد و کیفیت بالینی ماسک‌ها ادعا نمی‌شود.

\subsubsection{ارزیابی گروه‌آگاه در تصاویر پزشکی}

در داده‌ی پزشکی واحد استقلال غالباً شخص یا معاینه است، نه فایل. تصویرهای یک نوزاد می‌توانند نور، زمینه و شکل ویژه‌ی همان شبکیه را مشترک داشته باشند و ارزیابی گروه‌آگاه به مدل اجازه نمی‌دهد از این شباهت برای پیش‌بینی ظاهراً موفق موردی که قبلاً دیده است سود ببرد. این قاعده در تصویربرداری پزشکی استاندارد است و در این گزارش به‌عنوان شرط لازم ارزیابی به کار رفته، نه به‌عنوان نوآوری. در منبعی بدون شناسه‌ی بیمار، این قاعده فقط تا سطح شناسه‌ی در دسترس، یعنی معاینه، اعمال می‌شود.

\subsubsection{غربالگری پلاس در عمل}

تشخیص \lr{Plus} معمولاً در بخش مراقبت ویژه نوزادان با افتالموسکوپی غیرمستقیم انجام می‌شود و تأخیر یا کمبود متخصص خطر از دست رفتن پنجره‌ی درمان را بالا می‌برد. یک سامانه‌ی کمکی باید حساسیت بالا با اختصاصیت قابل قبول داشته باشد؛ در این پروژه آستانه‌ی یودن روی اعتبارسنجی این بده‌بستان را فقط به‌طور تقریبی متعادل می‌کند.

\subsubsection{سامانه‌های تحلیل تصویری پلاس}

سامانه‌های پیشین تحلیل تصویری پلاس بر اندازه‌گیری صریح رگ استوارند و پیچ‌خوردگی و گشادشدگی را از تصویر استخراج می‌کنند \cite{wittenberg2011,yildiz2020}. این خط پژوهشی به عملکرد بالایی رسیده است، پس فرض «بی‌ارزش بودن اندازه‌گیری صریح رگ» با شواهد سازگار نیست؛ پرسش این گزارش چیز دیگری است، آیا این اندازه‌ها روی یک بازنمایی آموخته‌شده اطلاعات مکمل اضافه می‌کنند یا نه. سه تفاوت میان این پروژه و آن سامانه‌ها تبیین‌کننده‌ی نتیجه‌ی این گزارش است؛ جدول ویژگی این پروژه ویژگی پهنای رگ یا گشادشدگی ندارد، شریان و ورید جدا نمی‌شوند، و ویژگی‌ها بر حسب پیکسل خام و روی کل تصویر محاسبه می‌شوند، نه با یک مقیاس آناتومیک در ناحیه‌ی قطب خلفی. کیفیت خود مرجع هم محدود است، چون توافق متخصصان بر تشخیص پلاس کامل نیست؛ این محدودیت سقف عملکرد هر مدل را تعیین می‌کند و دلیل دیگری است بر اینکه نتیجه‌ی این پروژه به‌عنوان پیش‌بینی برچسب ثبت‌شده در داده تفسیر شود، نه معادل تصمیم درمانی.

\subsubsection{جایگاه این کار}

این پژوهش تعریف بالینی \lr{Plus} را به یک مسئله‌ی رتبه‌بندی قابل اندازه‌گیری روی تصویر فوندوس تبدیل می‌کند، بدون ادعای جایگزینی معاینه. سهم افزایشی اندازه‌گیری صریح عروق در کنار یک بازنمایی تصویری یادگرفته‌شده و در کنار بازنمایی فضایی عروق سنجیده می‌شود. دستاورد این گزارش یک ادعای معماری نیست؛ یک پروتکل ارزیابی است که با قفل تقسیم، انتخاب مدل فقط با اعتبارسنجی، استنباط زوجی و سنجش انتقال به منبع کنارگذاشته اجرا شده است.
